\documentclass[11pt]{article}

\usepackage[margin=1in]{geometry}
\usepackage{amsmath,amssymb}
\usepackage{booktabs}
\usepackage{array}
\usepackage{hyperref}

\title{External Collapse-Ordering Invariance Assessment on Frozen HAOS-IIP Telemetry}
\author{Tomislav Rupi\'c \\ Independent Researcher}
\date{April 2026}

\begin{document}

\maketitle

\begin{abstract}
This paper freezes release \texttt{53.1} as a narrow post-foundations validation instrument layered above frozen HAOS-IIP telemetry. It is not a new phase, not a modification of HAOS-IIP core, and not a positive collapse-law claim. The instrument treats HAOS telemetry as a fixed black-box sensor, builds a score-ordered edge filtration on a graph-like artifact, perturbs that construction across locality, coupling, curvature, noise, and rewiring schedules, and asks only one bounded question: does the observed collapse ordering survive perturbation strongly enough to count as an invariant structural signal rather than an artifact of construction? On the current deterministic three-sector toy graph, a permissive toy control sensor yields \texttt{alpha $>$ beta $>$ gamma}, flip rate \(0.069\), persistent gap behavior, and an invariant-signal classification. The HAOS-IIP-backed adapter, which reuses only frozen telemetry, instead yields \texttt{alpha = beta $>$ gamma}, flip rate \(0.092\), vanishing gap behavior, and an artifact-of-construction classification. The bounded conclusion is therefore negative and useful: the external runner is now sufficient to separate protocol-level apparent signal from frozen-sensor structural support, but the present toy artifact does not yet earn a positive HAOS-IIP collapse-ordering claim.
\end{abstract}

\section{Scope}

This release sits closer to \texttt{45.1} and \texttt{50.3} than to a new phase paper.

\texttt{45.1} compressed frozen telemetry into a callable structural-stability oracle. \texttt{50.3} then defined the post-foundations discipline: validation should ask whether a candidate survives perturbation, purity checks, and scaling closure strongly enough to support bounded language. Release \texttt{53.1} applies that discipline to a new external question:
\[
\text{if a collapse ordering appears on a graph artifact, does it survive perturbation as structure?}
\]

This paper does not introduce:
\begin{itemize}
\item a new HAOS-IIP operator;
\item a new frozen telemetry channel;
\item a new sector claim;
\item a continuum or ontology claim;
\item a positive collapse-ordering law.
\end{itemize}

It does introduce:
\begin{itemize}
\item an external invariance-assessment protocol layered above frozen telemetry;
\item a strict distinction between protocol output and structural support;
\item a current negative result on the tested toy artifact under the HAOS-IIP-backed sensor.
\end{itemize}

\section{External-Layer Rule}

The architecture is intentionally separated:
\begin{itemize}
\item HAOS-IIP core remains frozen;
\item frozen telemetry remains the only authority sensor;
\item the validation runner is external and swappable;
\item the result is whatever survives perturbation under the frozen sensor.
\end{itemize}

The runner therefore treats HAOS metrics only as black-box observables. It does not reinterpret, extend, or reopen the frozen core.

\section{Protocol}

\subsection{Inputs}

Each run consumes:
\begin{itemize}
\item a graph \(G = (V,E)\) with sector labels on nodes;
\item edge attributes \(\mathrm{strength}(i,j)\) and distance \(d(i,j)\);
\item locality parameter \(\lambda\);
\item sector-coupling weights;
\item a curvature-penalty schedule;
\item a black-box HAOS sensor returning
\[
\{\Delta \mathrm{persistence},\; k_{\ast},\; \mathrm{safety\_margin},\; \mathrm{confidence}\}
\]
per inclusion step.
\end{itemize}

\subsection{Filtration}

Edges are scored by
\[
\mathrm{score}(i,j)
=
\mathrm{strength}(i,j)\,
\exp\!\left(-\frac{d(i,j)}{\lambda}\right)\,
\mathrm{curvature\_penalty}(i,j).
\]

The runner sorts edges by descending score and incrementally includes them to form a filtration \(G_t\). After each inclusion step it calls the black-box sensor and records the resulting time series.

\subsection{Sweeps and perturbations}

The current deterministic schedule contains \(36\) runs across:
\begin{itemize}
\item locality schedules with coarse, baseline, and fine \(\lambda\) steps;
\item sector-coupling rescalings;
\item curvature-penalty rescalings;
\item weight noise in the \(\pm 5\%\) to \(\pm 10\%\) range;
\item bounded random reshuffling on roughly \(10\%\) of the edges.
\end{itemize}

\subsection{Strict outputs}

The runner emits only:
\begin{itemize}
\item collapse ordering by sector;
\item consistency across trials, reported as pairwise flip rate;
\item gap behavior, reported as \texttt{persists}, \texttt{shrinks}, or \texttt{vanishes};
\item one bounded assessment: invariant structural signal or artifact of construction.
\end{itemize}

No interpretation beyond invariance assessment is part of the output contract.

\section{Current Deterministic Artifact}

The tested artifact is a small three-sector graph with sectors \texttt{alpha}, \texttt{beta}, and \texttt{gamma}. Its role is deliberately modest:
\begin{itemize}
\item it is large enough to produce nontrivial filtration orderings;
\item it is small enough for deterministic perturbation scans;
\item it is not treated as a physical model.
\end{itemize}

Two sensor backends were applied to the same artifact:
\begin{itemize}
\item a toy control sensor, used only as a protocol sanity check;
\item a HAOS-IIP-backed sensor adapter, built only from frozen telemetry.
\end{itemize}

The toy sensor is not authoritative. Its job is to show that the runner can register a robust ordering when the backend is permissive. The HAOS-IIP-backed adapter is the meaningful read because it uses the frozen telemetry layer already accepted by the repository.

\section{Frozen Quantitative Summary}

\begin{table}[h!]
\centering
\small
\begin{tabular}{@{}p{2.5cm}p{2.7cm}p{2.0cm}p{1.5cm}p{1.8cm}p{3.4cm}@{}}
\toprule
backend & collapse ordering & \(k_{\ast}\) steps & flip rate & gap behavior & assessment \\
\midrule
toy control &
\texttt{alpha $>$ beta $>$ gamma} &
\texttt{11, 12, 15} &
\(0.069\) &
\texttt{persists} &
\texttt{invariant structural signal} \\
HAOS-IIP-backed &
\texttt{alpha = beta $>$ gamma} &
\texttt{10, 10, 12} &
\(0.092\) &
\texttt{vanishes} &
\texttt{artifact of construction} \\
\bottomrule
\end{tabular}
\caption{Current deterministic invariance-assessment read on the three-sector toy artifact.}
\end{table}

The important row is the second one. Under the frozen telemetry adapter, the apparent ordering does not survive as a clean ordered ladder. The first two sectors hit \(k_{\ast}\) simultaneously, and the adjacent gap does not survive perturbation strongly enough to support a structural-ordering claim.

\section{What 53.1 Now Supports}

The strongest honest statement supported by this release is:

\begin{quote}
frozen HAOS-IIP telemetry can now be used as a black-box invariance-assessment sensor inside an external validation runner, and on the present toy artifact that runner rejects a positive collapse-ordering claim.
\end{quote}

That is not a failure of the instrument. It is the first useful property the instrument should have. A validation runner is more valuable when it can kill an apparent signal honestly than when it can always recover one.

\section{What 53.1 Refuses to Say}

This release does not justify any of the following:
\begin{itemize}
\item that collapse ordering is already a validated HAOS-IIP observable;
\item that \texttt{alpha $>$ beta $>$ gamma} is a structural law;
\item that the current toy graph is more than a probe artifact;
\item that the toy control sensor should be treated as evidence;
\item that any geometry-like, force-like, or continuum-like interpretation is now earned.
\end{itemize}

The negative HAOS-IIP-backed outcome is therefore part of the result, not an obstacle to it.

\section{Why This Matters}

Without this runner, one could easily mistake a construction-dependent ordering for a structural signal. The present release closes that loophole in a bounded way:
\begin{itemize}
\item the core remains frozen;
\item the probe layer is external;
\item the backend can be swapped;
\item only invariant survival under the frozen backend counts.
\end{itemize}

That is exactly the kind of discipline a post-foundations validation program should add before stronger language is allowed.

\section{Right Next Move}

The next practical step is not to upgrade the claim. It is to improve the artifact class:
\begin{itemize}
\item test graph families beyond the present toy artifact;
\item introduce sector layouts whose ties are not built into the construction;
\item keep the same runner and the same frozen sensor interface;
\item accept a positive statement only if the HAOS-IIP-backed path, not the toy control path, preserves the gap.
\end{itemize}

In other words, the current runner should now be reused, not reinterpreted.

\section{Conclusion}

Release \texttt{53.1} freezes an external collapse-ordering invariance-assessment protocol on frozen HAOS-IIP telemetry. The protocol is intentionally narrow, black-box, and perturbation-first. On the current deterministic toy artifact, the permissive control backend produces an apparent ordered signal, but the HAOS-IIP-backed backend classifies the same artifact as construction-sensitive because the leading gap vanishes under perturbation and the first two sectors collapse together. The bounded conclusion is therefore methodological and negative: the validation runner is now useful, but no positive HAOS-IIP collapse-ordering structural signal is yet earned from the present artifact.

\end{document}
